%/*% pdflatex -file-line-error -interaction nonstopmode %

\documentclass{article}
\usepackage{graphicx}
\usepackage{amsmath}
\usepackage{amsfonts}
\usepackage{hyperref}
\usepackage{exercises}
\usepackage{amsthm}
\usepackage{yhmath}
\newcommand*{\oldneg}{\mathord{\sim}}

\theoremstyle{definition}
\newtheorem{definition}{Definition}

\setlength{\parindent}{.5in}
\setlength{\parskip}{12pt}

\title{1.4 Exercises}
\begin{document}
\maketitle

\section*{1}
$4$

\section*{2}
If $a$ and $c$ are poth perfect squares, we can check to see if $b=2mn$ for $a=m^2,\ c=n^2$. If this is the case, we can use the perfect square trinomial 
pattern. If we only have two terms $a$ and $c$, we could then possibly use the difference of squares pattern.

\section*{3}
$2x^2+5x+3\newline\rightarrow (2x+3)(x+1)$

\section*{4}
$3n^2+7n+4\newline\rightarrow (3n+4)(n+1)$

\section*{5}
$4r^2+5r+1\newline\rightarrow (4r+1)(r+1)$

\section*{6}
$6p^2+5p+1\newline\rightarrow (3p+1)(2p+1)$

\section*{7}
$11z^2+2z-9\newline\rightarrow (11z-9)(z+1)$

\section*{8}
$15x^2-2x-8\newline\rightarrow (5x-4)(3x+2)$

\section*{9}
Cannot be factored.

\section*{10}
$14m^2+m-3\newline\rightarrow (7m-3)(2m+1)$

\section*{11}
$9d^2-13d-10\newline\rightarrow (9d+5)(d-2)$

\section*{12}
$5x^2+14x-3\newline\rightarrow (5x-1)(x+3)$

\section*{13}
$9x^2-1\newline\rightarrow (3x)^2-1\newline\rightarrow (3x+1)(3x-1)$

\section*{14}
$4r^2-25\newline\rightarrow (2r)^2-25\newline\rightarrow (2r+5)(2r-5)$

\section*{15}
$49n^2-16\newline\rightarrow (7n)^2-16\newline\rightarrow (7n+4)(7n-4)$

\section*{16}
$16s^2+8s+1\newline\rightarrow (4s)^2+8s+1\newline\rightarrow (4s+1)^2$

\section*{17}
$49x^2+70x+25\newline\rightarrow (7x)^2+70x+5^2\newline\rightarrow (7x+5)^2$

\section*{18}
$64w^2+144w+81\newline\rightarrow (8w)^2+144w+9^2\newline\rightarrow (8w+9)^2$

\section*{19}
$9p^2-12p+4\newline\rightarrow (3p)^2-12p+2^2\newline\rightarrow (3p-2)^2$

\section*{20}
$25t^2-30t+9\newline\rightarrow (5t)^2-30t+3^2\newline\rightarrow (5t-3)^2$

\section*{21}
$36x^2-84x+49\newline\rightarrow (6x)^2-84x+7^2\newline\rightarrow (6x-7)^2$

\section*{22}
$12x^2-4x-40\newline\rightarrow 4(3x^2-x-10)\newline\rightarrow 4((3x+5)(x-2))$

\section*{23}
$18z^2+36z+16\newline\rightarrow 2(9z^2+18z+8)\newline\rightarrow 2((3z+2)(3z+4))$

\section*{24}
$32v^2-2\newline\rightarrow 2(16v^2-1)\newline\rightarrow 2((4v+1)(4v-1))$

\section*{25}
$6u^2-24u\newline\rightarrow 6u(u-4)$

\section*{26}
$12m^2-36m+27\newline\rightarrow 3(4m^2-12m+9)\newline\rightarrow 3((2m)^2-12m+3^2)\newline\rightarrow 3((2m-3)^2)$

\section*{27}
$20x^2+124x+24\newline\rightarrow 4(5x^2+31x+6)\newline\rightarrow 4((5x+1)(x+6))$

\section*{28}
$21x^2-77x-28\newline\rightarrow 7(3x^2-11x-4)\newline\rightarrow 7((3x+1)(x-4))$

\section*{29}
$-36n^2+48n-15\newline\rightarrow -3(12n^2-16n+5)\newline\rightarrow -3((6n-5)(2n-1))$

\section*{30}
$-8y^2+28y-60\newline\rightarrow -4(2y^2-7y+15)\newline\rightarrow -4(2y^2-7y+15)$

\section*{31}
When the monomial was factored out, it was not taken out from the term $-36$. Instead, it should have been:\newline$4x^2-36\newline\rightarrow 
4(x^2-9)\newline\rightarrow 4((x+3)(x-3))$

\section*{32}
$16x^2-1=0\newline\rightarrow (4x+1)(4x-1)=0\newline\rightarrow x=\frac{1}{4}\vee x=-\frac{1}{4}$

\section*{33}
$11q^2-44=0\newline\rightarrow 11(q^2-4)=0\newline\rightarrow 11(q+2)(q-2)=0\newline\rightarrow q=2\vee q=-2$

\section*{34}
$14s^2-21s=0\newline\rightarrow (7s-0)(2s-3)=0\newline\rightarrow s=0\vee s=\frac{3}{2}$

\section*{35}
$45n^2+10n=0\newline\rightarrow (5n-0)(9n+2)=0\newline\rightarrow n=0\vee n=-\frac{2}{9}$

\section*{36}
$4x^2-20x+25=0\newline\rightarrow (2x)^2-20x+5^2=0\newline\rightarrow (2x-5)^2=0\newline\rightarrow x=\frac{5}{2}$

\section*{37}
$4p^2+12p+9=0\newline\rightarrow (2p)^2+12p+3^2=0\newline\rightarrow (2p+3)^2=0\newline\rightarrow p=-\frac{3}{2}$

\section*{38}
$15x^2+7x-2=0\newline\rightarrow (5x-1)(3x+2)=0\newline\rightarrow x=\frac{1}{5}\vee x=-\frac{2}{3}$

\section*{39}
$6r^2-7r-5=0\newline\rightarrow (2r+1)(3r-5)=0\newline\rightarrow x=-\frac{1}{2}\vee x=\frac{5}{3}$

\section*{40}
$36z^2+96z+15=0\newline\rightarrow 3(12z^2+32z+5)=0\newline\rightarrow 3(6z+1)(2z+5)=0\newline\rightarrow z=-\frac{1}{6}\vee z=-\frac{5}{2}$

\section*{41}
$4x^2-19x-5=0\newline\rightarrow (4x+1)(x-5)=0\newline\rightarrow x=-\frac{1}{4}\vee x=5$

\section*{42}
$3x^2-8x+5=0\newline\rightarrow (3x-5)(x-1)=0\newline\rightarrow x=\frac{5}{3}\vee x=1$

\section*{43}
$5x^2-27x-18=0\newline\rightarrow (5x+3)(x-6)=0\newline\rightarrow x=-\frac{3}{5}\vee x=6$

\section*{44}
$3x^2-3x=0\newline\rightarrow 3x(x-1)=0\newline\rightarrow (3x-0)(x-1)=0\newline\rightarrow x=0\vee x=1$ 

\section*{45}
$11x^2-19x-6=0\newline\rightarrow (11x+3)(x-2)=0\newline\rightarrow x=-\frac{3}{11}\vee x=2$

\section*{46}
$16x^2-2x-5=0\newline\rightarrow (8x-5)(2x+1)=0\newline\rightarrow x=\frac{5}{8}\vee x=-\frac{1}{2}$

\section*{47}
$15x^2-5x-20=0\newline\rightarrow (15x-20)(x+1)=0\newline\rightarrow x=\frac{4}{3}\vee x=-1$

\section*{48}
$18x^2-6x-4=0\newline\rightarrow (6x-4)(3x+1)=0\newline\rightarrow x=\frac{2}{3}\vee x=-\frac{1}{3}$

\section*{49}
$12x^2+5x-7=0\newline\rightarrow (12x-7)(x+1)=0\newline\rightarrow x=\frac{7}{12}\vee x=-1$

\section*{50}
$(2x)^2=36\newline\rightarrow (2x)^2-36=0\newline\rightarrow (2x+6)(2x-6)=0\newline\rightarrow x=-3\vee x=3$, and thus, because only 
positive values are accepted, $x=3$.

\section*{51}
$(3x+1)x=30\newline\rightarrow 3x^2+x-30=0\newline\rightarrow (3x+10)(x-3)=0\newline\rightarrow x=-\frac{10}{3}\vee x=3$, and thus, because only positive 
values are accepted, $x=3$.

\section*{52}
$\frac{1}{2}2x(5x-2)=115\newline\rightarrow 5x^2-2x-115=0\newline\rightarrow (5x+23)(x-5)=0\newline\rightarrow x=-\frac{23}{5}\vee x=5$, and thus, 
because only positive values are accepted, $x=5$. 

\section*{53}
$2x^2-4x-8=-x^2+x\newline\rightarrow 3x^2-5x-8=0\newline\rightarrow (3x-8)(x+1)=0\newline\rightarrow x=\frac{8}{3}\vee x=-1$

\section*{54}
$24x^2+8x+2=5-6x\newline\rightarrow 24x^2+14x-3=0\newline\rightarrow (6x-1)(4x+3)=0\newline\rightarrow x=\frac{1}{6}\vee x=-\frac{3}{4}$

\section*{55}
$18x^2-22x=28\newline\rightarrow 18x^2-22x-28=0\newline\rightarrow (9x+7)(2x-4)=0\newline\rightarrow x=-\frac{7}{9}\vee x=2$

\section*{56}
$13x^2+21x=-5x^2+22\newline\rightarrow 18x^2+21x-22=0\newline\rightarrow (6x+11)(3x-2)=0\newline\rightarrow x=-\frac{11}{6}\vee x=\frac{2}{3}$

\section*{57}
$x=4x^2-15x\newline\rightarrow 4x^2-16x=0\newline\rightarrow (4x-0)(x-4)=0\newline\rightarrow x=0\vee x=4$

\section*{58}
$(x+8)^2=16-x^2+9x\newline\rightarrow x^2+16x+64=16-x^2+9x\newline\rightarrow 2x^2+7x+48=0\newline\rightarrow$, because this cannot be factored, 
the equation has no solution. 

\section*{59}
$2x^3-5x^2+3x\newline\rightarrow x(2x^2-5x+3)\newline\rightarrow x(2x+1)(x-3)$

\section*{60}
$8x^4-8x^3-6x^2\newline\rightarrow 2x^2(4x^2-4x-3)\newline\rightarrow 2x^2(2x+1)(2x-3)$

\section*{61}
$9x^3-4x\newline\rightarrow x(9x^2-4)\newline\rightarrow x(3x+2)(3x-2)$

\section*{62}
$4=(2+2x)(1+2x)-2\newline\rightarrow 4=4x^2+6x\newline\rightarrow 4x^2+6x-4=0\newline\rightarrow 2(2x^2+3x-2)=0\newline\rightarrow 2(2x-1)(x+2)=0\newline
\rightarrow x=\frac{1}{2}\vee x=-2$, and thus, because only positive values are accepted, $x=\frac{1}{2}$

\section*{63}
$96=(12+2x)(8+2x)-96\newline\rightarrow 4x^2+40x-96=0\newline\rightarrow (4x+48)(x-2)=0\newline\rightarrow x=-12\vee x=2$ Because only positive values 
are accepted, the width of the border should be $2$ ft. 

\section*{64}
$R(x)=(500-20x)(45+5x)\newline\rightarrow -100(x-25)(x+9)$. Thus, the zeros of the function are $25, -9$, meaning the vertex lies on the line $x=-\frac
{25-9}{2}=-8$, and thus, the store should sell each board for $-20(-8)=160$ more, and, therefore, $D$ is the correct answer.

\section*{65}
First, we set up a quadratic equation to model this situation:\newline $R(x)=(330+15x)\times (6-0.25x)$\newline Next, we find the vertex of this equation, 
by finding it's zeros:\newline $-3.75(x+22)(x-24)$, and thus, the zeros are $-22$ and $24$, and thus the $x$-coordinate of the vertex is 
$\frac{-22+24}{2}=1$, and thus, the maximum value is $(330+15(1))(6-0.25(1))=\$1983.75$ when each sandwich is sold for $6-0.25(1)=\$5.75$

\section*{66}
$714=(25+4x)(21+2x)-525\newline\rightarrow 8x^2+134x-714=0\newline\rightarrow (4x-17)(2x+42)=0\newline\rightarrow x=\frac{17}{4}\vee x=-21$, and thus, 
because only positive values are accepted, the width of the mat is $42$ ft., while the width at the top and bottom of the paintings is $29.5$ ft.

\section*{67}
(a): The girth is $108-36=72$ ft.\newline
(b): For the width: $72=2w+2h\newline\rightarrow w=36-h$\newline For the volume: $V=wlh=36(36-h)h$\newline
(c): We can rewrite the above formula as a quadratic formula, in order to find the $x$-coordiante of its vertex:\newline
$V=-36h^2+1296h\newline\rightarrow x cord =-\frac{1296}{2(-36)}=18$, and thus, the maximum volume is $-36(18)^2+1296(18)=11664$. This is the 
maximum volume, as it is the maximum height of the quadratic function. So, a height of $18$ in. and a width of $36-18=18$ in. maximize the 
volume of the box, at $11664 in^3$.

\section*{68}
$2x(5x-4)=(3x+2)(x+1)\newline\rightarrow 10x^2-8x=3x^2+5x+2\newline\rightarrow 7x^2-13x-2=0\newline\rightarrow (7x+1)(x-2)=0\newline\rightarrow 
x=-\frac{1}{7}\vee x=2$, and thus, because only positive values are accpeted, $x=2$. 
\end{document}

